
\documentclass[11pt,a4paper]{article}
\usepackage[a4paper,margin=2.05cm]{geometry}
\usepackage[spanish,es-noquoting]{babel}
\usepackage[T1]{fontenc}
\usepackage[utf8]{inputenc}
\usepackage{lmodern,microtype}
\usepackage{amsmath,amssymb,amsthm,mathtools}
\usepackage{booktabs,longtable,array,enumitem}
\usepackage{xcolor,hyperref}
\pagecolor{white}\color{black}
\setlength{\parindent}{0pt}
\setlength{\parskip}{0.65em}
\newtheorem{teorema}{Teorema}
\newtheorem{observacion}[teorema]{Observación}
\title{\bfseries N64 --- Ley de neutralidad Witt--dual para frontera saturada\\
\large Justificación de \((BA_W)\mathbf1=0\), selección de \(r=102\) y canonización del cierre}
\author{HMT / auditoría técnica}
\date{}
\begin{document}
\maketitle
\begin{abstract}
N64 cierra la justificación conceptual que quedaba amarilla en N63. En una frontera saturada, el carry uniforme \(c=111111\) define la unidad visible \(\mathbf1\). Como \(A_W^2=-I\), la matriz de Witt/Golay transporta esa unidad a una unidad dual
\[
\omega=A_W\mathbf1=211111.
\]
La ley de cierre saturado no es neutralidad visible \(B\mathbf1=0\), sino neutralidad Witt--dual:
\[
\boxed{(BA_W)\mathbf1=B\omega=0.}
\]
Con \(q,a,c,\operatorname{colw}\) y la poda decimal, esta ley selecciona de forma única la matriz real y por tanto deriva \(r=B\mathbf1=102\). Sin poda decimal deja el par \(\pi/e\), que la poda rompe. Se aclara además que la condición no es la única forma lineal imaginable; es la forma canónica inducida por la estructura compleja ternaria de Witt.
\end{abstract}
\section{Unidad visible y unidad dual}
La matriz de Witt satisface:
\[
A_W^2=-I.
\]
Sea:
\[
\mathbf1=(1,1,1,1,1,1)^T.
\]
Entonces:
\[
\boxed{\omega=A_W\mathbf1=211111.}
\]
Además:
\[
A_W\omega=-\mathbf1.
\]
Así, \(\omega\) es la unidad dual de Witt asociada a la unidad visible.
\section{Ley de frontera saturada}
En la sexta celda se tiene:
\[
c=111111.
\]
Esto expresa saturación completa de carry. La ley de cierre saturado exige que la frontera no tenga carga neta en la coordenada dual de Witt:
\[
\boxed{(BA_W)\mathbf1=0.}
\]
Como:
\[
(BA_W)\mathbf1=B(A_W\mathbf1)=B\omega,
\]
la condición equivalente es:
\[
\boxed{B\omega=0.}
\]
\section{Visible no es dual}
La neutralidad visible sería:
\[
B\mathbf1=0.
\]
Pero la matriz real tiene:
\[
B\mathbf1=102.
\]
Por tanto la frontera no es visible-neutral. Es dual-neutral:
\[
(BA_W)\mathbf1=000.
\]
La carga visible \(r=102\) es precisamente la sombra visible de una frontera cerrada en la coordenada dual.
\section{Selección en el cilindro decimal}
Con las condiciones derivadas:
\[
q=012120,\quad a=111101,\quad c=111111,\quad \operatorname{colw}=223232,
\]
y con la poda decimal de la quinta triada, quedan cuatro matrices. Al imponer:
\[
B\omega=0,
\]
queda una única matriz:
\[
\boxed{
B=
\begin{pmatrix}
222220\\
021222\\
102011
\end{pmatrix}.
}
\]
Entonces:
\[
\boxed{r=B\mathbf1=102.}
\]
\section{Sin poda decimal}
Sin poda decimal, las condiciones \(q,a,c,\operatorname{colw}\) dejan ocho matrices. La neutralidad dual deja dos:
\[
(021222,222220,102011),
\]
\[
(222220,021222,102011).
\]
Son el par intercambiado \(\pi/e\). La poda \(729/1000\) selecciona la orientación real:
\[
(222220,021222,102011).
\]
\section{Canon frente a arbitrariedad}
Hay otras formas lineales arbitrarias que podrían separar candidatos. N64 no afirma unicidad entre todas las formas lineales. La afirmación correcta es:
\[
\boxed{B\omega=0\text{ es el selector canónico inducido por }A_W^2=-I.}
\]
Dentro de la órbita de la unidad visible bajo \(A_W\):
\[
\mathbf1,\quad A_W\mathbf1=\omega,\quad -\mathbf1,\quad -\omega,
\]
la neutralidad visible no sirve y la neutralidad dual selecciona el cierre correcto.
\section{Dictamen}
\[
\boxed{\text{verde: la ley saturada es neutralidad Witt--dual }B(A_W\mathbf1)=0.}
\]
\[
\boxed{\text{verde: }r=102\text{ se deriva como carga visible inducida.}}
\]
\[
\boxed{\text{verde: sin poda queda el par }\pi/e;\text{ con poda queda la orientación real.}}
\]
\[
\boxed{\text{verde disciplinado: canon, no unicidad entre formas lineales arbitrarias.}}
\]
\end{document}
